\section{Limitations}
\label{sec:limitations}

\paragraph{Fault class.}
The peer-to-peer transport of the evaluated platform is encrypted, so
network-level injection cannot be made selective with respect to consensus
message type. Experiment X5 therefore realizes \emph{omission and duplication}
faults (Definition~\ref{def:faults}) and not equivocation or signature forgery.
The measured $\pd$ and $\pf$ characterize the detector against this fault class
only; a detector facing equivocating validators would require protocol-level
instrumentation and is left for future work. All statements in
Section~\ref{sec:safety} are conditional on this fault class.

\paragraph{Absolute throughput.}
Under the \RNM{} protocol each validator receives a small fixed CPU quota, so
absolute throughput figures are far below what dedicated hardware would deliver
and are not comparable with published benchmarks. Every claim in this paper
concerns the exponents $\beta$ and $\gamma$, the coverage of prediction
intervals, and the sizing decisions derived from them. Statement~\ref{st:qinvariance}
is what licenses this separation, and it is tested rather than assumed.

\paragraph{Measured range.}
The largest directly measured validator set is bounded by host memory.
Extrapolation beyond it rests on the calibrated model and on the
simulation of experiment X6; the prediction intervals of
Section~\ref{sec:pi} quantify, but do not eliminate, the associated risk.
Coverage was validated only on held-out counts inside the reachable range.

\paragraph{Single host and single platform.}
All measurements come from one host and one consensus implementation. Wide-area
behaviour is emulated by kernel-level latency injection rather than observed on
a real multi-region deployment, and hardware heterogeneity is not represented.

The two contributions are not equally exposed to this.
Theorem~\ref{thm:identifiability} is a statement about least squares, and its
validity does not depend on the platform at all; what depends on the platform
are the values of $\beta$ and $\gamma$. This yields a prediction that a second
implementation can falsify cheaply: fitting $(\beta,\gamma)$ on another
QBFT-family client and then re-running the path re-sampling of
Section~\ref{sec:experiments} should reproduce the line $\ahat=\gamma\lambda-\beta$
with \emph{that} client's coefficients, and in particular should place the sign
change at $\lambda=\beta/\gamma$. We report this replication in
experiment~X9 for one additional client; a failure there would undermine the
model~\eqref{eq:bifactor} rather than the theorem. The
\RNM{} protocol is more exposed, since Assumption~\ref{as:separable} concerns
the scheduler and the runtime, and a client with a different threading model may
violate it at quotas where ours does not.

\paragraph{Asymmetry between theory and measurement.}
We state plainly that the analytical part of this work is more complete than the
empirical part. The theorems hold as stated; the calibrated coefficients rest on
a campaign bounded by what a four-core runner can host, and
Assumption~\ref{as:separable} in particular is supported over a narrow quota
range. Readers should treat the numerical values as provisional and the
structural claims --- the form of the bias, the shape of the feasible set --- as
the durable part.

\paragraph{Model assumptions.}
The separability of Assumption~\ref{as:separable}, the exchangeable correlation
of Assumption~\ref{as:rho}, and the approximate log-normality underlying
Statement~\ref{st:pi} are approximations. Each is stated explicitly and each is
accompanied by a test in Section~\ref{sec:results}; where a test rejects an
assumption, the corresponding claim is restricted rather than retained.

\paragraph{Cost model.}
Theorem~\ref{thm:window} assumes only that $\Cost$ is nondecreasing in the
validator count. Deployments where additional validators carry
non-monotone institutional or governance cost fall outside this assumption,
although the feasibility window itself remains valid.
